% !TEX root = ../main.tex
\chapter{The Four Prices of Money}\label{ch:fourprices}
\begin{paracol}{2}

\begin{sources}
\begin{tabularx}{\linewidth}{@{}l l L@{}}
\toprule
\textbf{Segment} & \textbf{Length} & \textbf{Content}\\
\midrule
\seg{1}{1}{L1.1 The Big Picture} & 19:26 & Where the course comes from: Stigum, the \emph{Financial Times}, the history of monetary and financial thought; the ``money view''; dialogue between central banking and modern finance.\\
\seg{1}{2}{L1.2 Prerequisites?} & 7:22 & IS/LM and the Fisher diagram: the two textbook ways of teaching money and banking, and why this course starts elsewhere.\\
\seg{1}{3}{L1.3 What is a Bank, a Shadow Bank, a Central Bank} & 12:10 & Balance sheets as the analytical tool; solvency versus liquidity; shadow banks; the Fed's balance sheet in 2012.\\
\seg{1}{4}{L1.4 Central Themes} & 13:09 & Plan of the course; the four prices of money; Bagehot's \emph{Lombard Street}; ``all banking is a swap of IOUs''.\\
\seg{1}{5}{L1.5 Reading: Allyn Young} & 3:10 & Presentation of the first reading.\\
\bottomrule
\end{tabularx}

\smallskip
\textbf{Files.} Mehrling's lecture notes \file{money-lecture_notes-Lec 01--The Four Prices of Money.pdf};
captions \file{materials/week1/srt/1 - *.srt}; transcripts \file{materials/transcripts/L01-P0*.txt}.
\textbf{Reading.} Allyn Young, four chapters from \emph{The Book of Popular Science} (1924), file
\file{Allyn Young.pdf}, with the study questions \file{Allyn Young(study).pdf}; discussed in
\cref{sec:fp-young} and again in lecture~3.
\end{sources}

\section*{Motivation}
The first lecture is an orientation. It says where the course stands among the ways money and banking
are usually taught, introduces the one analytical tool that will be used everywhere (the balance sheet),
and names the four prices that the whole course tries to explain:
\begin{itemize}
  \item \textbf{par}, the price of one kind of money in terms of another kind today;
  \item \textbf{interest}, the price of money today in terms of money tomorrow;
  \item the \textbf{exchange rate}, the price of domestic money in terms of foreign money;
  \item the \textbf{price level}, the price of money in terms of goods.
\end{itemize}
The thread of the lecture is that the standard textbooks describe the banking system of 1950, while
the system of today is \emph{global} and \emph{market-based}, and the problems it runs into are
problems of \emph{liquidity}, not only of solvency. For a physicist, the useful attitude is that of an
experimentalist entering a new lab: the course does not start from a model and apply it, but starts
from what practitioners actually do and tries to extract the regularities \ts{1}{2}{4:41}.

% ------------------------------------------------------------------
\section{The big picture}\label{sec:fp-bigpicture}

\paragraph{Where the knowledge comes from.} Mehrling has never traded in the money markets: his
knowledge is ``book learning'', but from practitioners' books, above all Marcia Stigum's
\emph{The Money Market}, and from the \emph{Financial Times}, rather than from academic journals
\ts{1}{1}{0:29}. His academic work is in the history of monetary and financial economics, in two
books: \emph{The Money Interest and the Public Interest}, a history of American monetary thought, and
\emph{Fischer Black and the Revolutionary Idea of Finance}, a history of the rise of modern finance
\ts{1}{1}{0:50}. These are the intellectual foundations of the course, and they explain why the
readings include authors such as Allyn Young, Walter Bagehot, Ralph Hawtrey and Edward Shaw
\ts{1}{1}{1:10}.\begin{aside}[Mehrling's books]
Perry Mehrling (Barnard College, Columbia University, at the time of the lectures; later Boston
University) wrote \emph{The Money Interest and the Public Interest: American Monetary Thought,
1920--1970} (Harvard University Press, 1997), \emph{Fischer Black and the Revolutionary Idea of
Finance} (Wiley, 2005) and \emph{The New Lombard Street: How the Fed Became the Dealer of Last Resort}
(Princeton University Press, 2011), the book that accompanies this course.
\end{aside}

\paragraph{A course about the present.} Although its foundations are historical, this is not a course
in the history of economic thought: it uses that history as a frame to understand the world as it is
now \ts{1}{1}{1:31}. The big fact to understand is \emph{financial globalization} (or ``global
financialization''). The lesson of the history of monetary thought is that every generation has to
work out for itself the world in front of it: some ideas carry over from the past, but the system is
not constant, and today it is ``a strange new world'' \ts{1}{1}{1:51}. This is why Mehrling does not use
the standard money and banking textbooks: they evolved from the textbooks of 1950, and a bank in 1950
was a different thing from a bank today \ts{1}{1}{2:53}. Better to put them in the library and start
again from the real world, accepting that the result is a little ``looser'' but attached to reality
\ts{1}{1}{3:14}.

\paragraph{Listening to the bankers.} Being attached to reality means listening to practitioners and
trying to understand their world \ts{1}{1}{3:35}. Practitioners know their own piece very well (repo,
Eurodollars, \dots) but have no picture of how the system as a whole works; an academic can add value
by noticing the patterns and the connections between the pieces \ts{1}{1}{3:55}. This is not easy:
there will be no ``take the first derivative and set it equal to zero'' in this course \ts{1}{1}{4:36}.

\begin{intuition}[Reading Stigum as an anthropologist]
Stigum's \emph{The Money Market} is the desk reference of people who trade the money markets, ``the
bible'' of the field \ts{1}{1}{4:56}. Mehrling's advice is to read it as an anthropologist reads the
sacred text of a foreign tribe \ts{1}{1}{5:58}: not to learn to trade, but to understand how the
members of the tribe think about their world. Read for every detail, it ``can blow your mind''; read for
the way of thinking, it is a window into a rather special world \ts{1}{1}{6:19}.
\end{intuition}\begin{aside}[Stigum]
Marcia Stigum (1934--2005) first published \emph{The Money Market: Myth, Reality, and Practice} in 1978;
the fourth edition (with Anthony Crescenzi, McGraw-Hill, 2007) is the one mentioned in class, with a
new chapter on Eurodollars \ts{1}{1}{5:16}. Mehrling recalls that when he started teaching the course,
Stigum was the whole text, with no lecture notes.
\end{aside}

\paragraph{Reading the Financial Times.} In a sense the whole course is about learning to read the
\emph{Financial Times} (FT) and to ``translate it into English'' \ts{1}{1}{7:21}; the final exam contains
exactly that exercise, an article to be explained to one's roommate \ts{1}{1}{7:42}. Every lecture
after this one starts with an article of the day. Three examples from the paper of that day show the
kind of question the course is about \ts{1}{1}{8:23}:
\begin{enumerate}[(a)]
  \item German banks split over the role of the European Central Bank (ECB): should the ECB buy the
        government debt of Spain or Italy to support its price? What does it even mean for a central
        bank to buy the debt of a state \ts{1}{1}{8:43}?
  \item Chinese steel mills brace for a slowdown in demand. The monetary angle: in China, inventories of
        iron ore had been used as \emph{collateral} for borrowing; as the banks call in the loans, the
        borrowers dump the collateral, the beginning of a financial crisis \ts{1}{1}{9:04}.
  \item Foreign exchange brokers back plans of the trading platform EBS to rein in ``predatory'' traders:
        a battle between algorithmic high-frequency traders (``cyborgs'') and human dealers, who lose
        money to them \ts{1}{1}{9:45}.
\end{enumerate}

\subsection{The money view and its ancestors}
Asked which school of economics he belongs to (saltwater or freshwater, Keynesian or monetarist),
Mehrling answers that he identifies with the \emph{money view}, a name he invented himself
\ts{1}{1}{10:47}. It has two lines of intellectual ancestors \ts{1}{1}{11:08}, shown in
\cref{fig:fp-ancestors}.
\begin{itemize}
  \item \textbf{The American line.} It starts with Allyn Young, because central banking in the United
        States starts in his time (the Fed was founded in 1913), when there was no tradition of central
        banking thought to draw on \ts{1}{1}{11:49}. Then Alvin Hansen and Edward Shaw (the three
        generations of Mehrling's first book), and Hyman Minsky, who would have been the missing last
        chapter of that book \ts{1}{1}{12:09}.
  \item \textbf{The British line.} Britain had a central bank for a long time and was the centre of the
        world economy for a long time, so it has a long tradition of thinking about central banking
        \ts{1}{1}{12:50}: Walter Bagehot, Ralph Hawtrey, Richard Sayers (who taught at the London School of
        Economics, LSE) and Charles Goodhart \ts{1}{1}{13:10}. Almost all of them talked to bankers and
        central bankers: Goodhart spent most of his career at the Bank of England \ts{1}{1}{13:30}, and to
        teach at the LSE is to be a mile from the City of London, as teaching in Manhattan is to be a mile
        from Wall Street \ts{1}{1}{13:51}.
\end{itemize}
Some of them called themselves Keynesians (Hansen, in a way Minsky), others leaned towards monetarism
(Shaw, Hawtrey); Mehrling reads them as one continuous stream of thought that kept its eye on the real
world of its own time \ts{1}{1}{14:11}.

\begin{nfigure}
\begin{tikzpicture}[x=1cm,y=1cm,font=\small,
  p/.style={draw=defcol,fill=defcol!8,rounded corners=2pt,inner sep=3pt,align=center,minimum width=2.5cm},
  q/.style={draw=thmcol,fill=thmcol!8,rounded corners=2pt,inner sep=3pt,align=center,minimum width=2.5cm},
  arr/.style={-{Stealth[length=4pt]},semithick}]
  \node[font=\bfseries\small,defcol] at (0,0.8) {American line};
  \node[p] (y) at (0,0) {Allyn Young\\\scriptsize 1876--1929};
  \node[p] (h) at (0,-1.2) {Alvin Hansen\\\scriptsize 1887--1975};
  \node[p] (s) at (0,-2.4) {Edward Shaw\\\scriptsize 1908--1994};
  \node[p] (m) at (0,-3.6) {Hyman Minsky\\\scriptsize 1919--1996};
  \draw[arr] (y)--(h); \draw[arr] (h)--(s); \draw[arr] (s)--(m);
  \node[font=\bfseries\small,thmcol] at (8,0.8) {British line};
  \node[q] (b) at (8,0) {Walter Bagehot\\\scriptsize 1826--1877};
  \node[q] (ha) at (8,-1.2) {Ralph Hawtrey\\\scriptsize 1879--1975};
  \node[q] (sa) at (8,-2.4) {Richard Sayers\\\scriptsize 1908--1989};
  \node[q] (g) at (8,-3.6) {Charles Goodhart\\\scriptsize b.~1936};
  \draw[arr] (b)--(ha); \draw[arr] (ha)--(sa); \draw[arr] (sa)--(g);
  \node[draw=intucol,fill=intucol!10,very thick,rounded corners=3pt,align=center,inner sep=5pt] (mv) at (4,-1.8)
    {\textbf{the money view}\\\scriptsize (Mehrling)};
  \draw[arr,intucol,dashed] (y.east) to[out=0,in=150] (mv);
  \draw[arr,intucol,dashed] (m.east) to[out=0,in=-150] (mv);
  \draw[arr,intucol,dashed] (b.west) to[out=180,in=30] (mv);
  \draw[arr,intucol,dashed] (g.west) to[out=180,in=-30] (mv);
  \node[draw=fincol,fill=fincol!10,rounded corners=2pt,align=center,inner sep=3pt] (fb) at (4,-4.4)
    {Fischer Black, modern finance\\\scriptsize ``no separate theory of money''};
  \draw[{Stealth[length=4pt]}-{Stealth[length=4pt]},semithick,fincol] (fb) -- node[right,font=\scriptsize]{dialogue} (mv);
\end{tikzpicture}
\caption{The intellectual ancestors of the money view \ts{1}{1}{11:28}, and its opponent in the
dialogue that the course sets up \ts{1}{1}{16:14}.}\label{fig:fp-ancestors}
\end{nfigure}

\subsection{The challenge of modern finance}
Mehrling also learned modern finance, through his biography of Fischer Black. Finance holds that there
is no such field as monetary economics \ts{1}{1}{14:53}: in an efficient market every asset trades at its
efficient price, so there is nothing for liquidity to distort. In ``straight-ahead'' finance there is no
\emph{fire sale}: if a bank had to sell its loans, a price just a tiny bit below their true value would be
enough for someone to snatch them up, so prices never move far from value and there are no runs on banks
\ts{1}{1}{15:13}. The logical end of this view is that there is no reason for central banks to exist
\ts{1}{1}{15:54}.

The course is, in Mehrling's words, a \emph{dialogue} between the classical central banking perspective and
the modern finance perspective, two ideas that seem contradictory, rubbed against each other to make
sparks for new economic thinking \ts{1}{1}{16:14}. The opposition is not Keynesian versus monetarist
\ts{1}{1}{16:55}.

\paragraph{Where Mehrling stands.} Pressed by the student, he states the position he will defend for the
rest of the course \ts{1}{1}{17:38}: there \emph{is} a role for central banks in setting bounds on the
system, a modern version of the role that Bagehot had in mind in 1873 (``lend freely at a high rate against
good security'') \ts{1}{1}{17:59}. That rule does not quite work in the modern world, and the subtitle of his
book, \emph{How the Fed Became the Dealer of Last Resort}, signals the update. Financial markets do not work
perfectly; they run into problems, and banking is all about these problems, which are \emph{liquidity}
problems rather than problems of asymmetric information \ts{1}{1}{18:20}. The most compelling work on them
comes not from academic journals, ``fighting old battles'', but from the Bank of England and the Bank for
International Settlements (BIS), the central bankers' bank \ts{1}{1}{19:01}.\begin{history}[Bagehot's \emph{Lombard Street} (1873)]
Walter Bagehot (1826--1877), editor of \emph{The Economist}, published \emph{Lombard Street: A Description of
the Money Market} in 1873, after the panic that followed the failure of the discount house Overend, Gurney \&
Co.\ in 1866. His prescription for the Bank of England in a panic is usually summarised as ``lend freely, at a
high rate, against good collateral''; in his own words, loans ``should only be made at a very high rate of
interest'' and ``on all good banking securities, and as largely as the public ask for them''. Lombard Street,
in the City of London, takes its name from the Lombard (northern Italian) merchants and bankers who settled
there in the Middle Ages.\par
\emph{References:} W.~Bagehot, \emph{Lombard Street}, London: Henry S.~King, 1873, ch.~VII; Mehrling,
\emph{The New Lombard Street}, 2011, ch.~1.
\end{history}

% ------------------------------------------------------------------
\section{Prerequisites: two textbook approaches}\label{sec:fp-prereq}

The formal prerequisites are intermediate macroeconomics and intermediate microeconomics
\ts{1}{2}{0:08}. The course makes little use of them, but other money and banking courses start from them,
and at the end of the course the money view will be connected to both \ts{1}{2}{0:29}. It is worth seeing what
they are, because they are what the course is \emph{not} doing.

\paragraph{Intermediate macro: IS/LM.} The macro story determines national income $Y$ and the
\emph{nominal} interest rate $i$ at the crossing of two curves \ts{1}{2}{1:15}:
\begin{itemize}
  \item the IS curve (``investment = saving''): the pairs $(Y,i)$ at which the market for goods clears;
        a higher rate depresses investment, so the curve slopes down;
  \item the LM curve (``liquidity preference = money supply''): the pairs at which the demand for money
        equals the supply; higher income raises the demand for money, so the rate must rise to keep it
        equal to a given supply, and the curve slopes up.
\end{itemize}
Monetary policy is then a shift of the LM curve. A whole course can be built on the foundations of the LM
curve, the Taylor rule, or the argument that with inflation targeting the LM curve should be horizontal
\ts{1}{2}{1:35}. This course instead questions the very notion of a stable money demand curve and money supply
curve, at a more fundamental level \ts{1}{2}{1:55}.

\paragraph{Intermediate micro: the Fisher diagram.} The micro story determines the \emph{real} interest rate in a
two-period general equilibrium \ts{1}{2}{2:16}, shown in \cref{fig:fp-prereq}. With wealth $W$ in units of
today's goods, consumption $c_1$ today and $c_2$ tomorrow, and real rate $r$, the budget line is
\begin{equation}\label{eq:fp-budget}
  c_1 + \frac{c_2}{1+r} = W, \qquad \frac{\dd c_2}{\dd c_1} = -(1+r),
\end{equation}
so its slope is $-(1+r)$ in the $(c_1,c_2)$ plane (Mehrling's notes write $-1/(1+r)$, which is the slope with
the axes exchanged). Maximising utility $u(c_1,c_2)$ on this line gives the first-order condition
\[
  \frac{\partial u/\partial c_1}{\partial u/\partial c_2} = 1+r,
\]
the tangency between an indifference curve and the budget line \ts{1}{2}{2:36}. Two people with the same
wealth and the same production possibilities but different tastes choose different points: one consumes more
than she produces today and borrows, the other consumes less and lends \ts{1}{2}{2:58}. Financial markets allow
this inter-temporal trade, and the real interest rate is the price that clears it \ts{1}{2}{3:19}. The diagram
is named after Irving Fisher \ts{1}{2}{3:39}. A money and banking course built on it treats banks as
intermediaries between savers and investors, needed because of some imperfection such as asymmetric
information \ts{1}{2}{4:00}.

\begin{nfigure}
\begin{tikzpicture}[font=\small]
  \begin{scope}
    \draw[-{Stealth[length=4pt]}] (0,0) -- (4.2,0) node[below left]{income $Y$};
    \draw[-{Stealth[length=4pt]}] (0,0) -- (0,3.4) node[left]{$i$};
    \draw[thick,defcol] (0.5,3.0) -- (3.8,0.5) node[right]{IS};
    \draw[thick,thmcol] (0.6,0.6) -- (3.8,2.2) node[right]{LM};
    \fill (2.16,1.39) circle (1.5pt);
    \draw[dashed,thin] (2.16,0) -- (2.16,1.39) -- (0,1.39);
    \node[font=\bfseries] at (2,3.5) {macro: nominal rate};
  \end{scope}
  \begin{scope}[xshift=6.2cm,scale=0.95]
    % PPF: quarter circle of radius 3; production at angle 40 deg, P=(2.298,1.928), slope -cot 40 = -1.192
    % budget line through P; indifference curves c1^a c2^(1-a) tangent to it (a=0.35 lender, a=0.75 borrower)
    \draw[-{Stealth[length=4pt]}] (0,0) -- (4.6,0) node[below]{$c_1$};
    \draw[-{Stealth[length=4pt]}] (0,0) -- (0,5.0) node[left]{$c_2$};
    \draw[thick,intcol] (0,3) arc (90:0:3);
    \node[intcol,font=\scriptsize] at (0.55,2.75) {PPF};
    \draw[thick,fincol] (0,4.667) -- (3.916,0);
    \node[fincol,font=\scriptsize,anchor=west] at (2.05,3.45) {slope $-(1+r)$};
    \draw[fincol,thin] (2.0,3.42) -- (0.95,3.53);
    \draw[defcol,domain=0.78:3.6,smooth,variable=\x] plot (\x,{3.031*(1.371/\x)^0.5385});
    \draw[thmcol,domain=1.92:4.4,smooth,variable=\x] plot (\x,{1.166*(2.937/\x)^3});
    \fill (2.298,1.928) circle (2pt);
    \node[font=\scriptsize,anchor=east] at (2.25,1.75) {production};
    \fill[defcol] (1.371,3.031) circle (2pt);
    \fill[thmcol] (2.937,1.166) circle (2pt);
    \draw[dashed,thin,defcol] (1.371,0) -- (1.371,3.031);
    \draw[dashed,thin] (2.298,0) -- (2.298,1.928);
    \draw[dashed,thin,thmcol] (2.937,0) -- (2.937,1.166);
    \draw[decorate,decoration={brace,amplitude=3pt,mirror},defcol] (1.40,-0.08) -- (2.27,-0.08)
      node[midway,below=3pt,font=\scriptsize]{lend};
    \draw[decorate,decoration={brace,amplitude=3pt,mirror},thmcol] (2.33,-0.08) -- (2.91,-0.08)
      node[midway,below=3pt,font=\scriptsize]{borrow};
    \node[font=\bfseries\small] at (2.3,5.5) {micro: real rate};
  \end{scope}
\end{tikzpicture}
\caption{The two textbook starting points \ts{1}{2}{1:15}\ts{1}{2}{2:16}. Left: IS/LM fixes income and the nominal
interest rate. Right: the Fisher diagram; two agents produce at the same point of the production possibility
frontier (PPF) and trade along the budget line to the tangency with their own indifference curve: the one who
consumes less today lends, the other borrows.}\label{fig:fp-prereq}
\end{nfigure}

\paragraph{The strategy of this course.} The course does not start from an economic apparatus and ask how to
apply it to banking: it starts from the concrete realities of banking and tries to draw from them, inductively,
how the system works and what regularities it has \ts{1}{2}{4:41}. Mehrling warns that the first two weeks
feel disorienting, because there is nothing familiar to grab onto; it comes together later \ts{1}{2}{5:02}. His own
way of reading the FT is to ask what Bagehot, Minsky or Fischer Black, whose biographies he wrote, would make of
it \ts{1}{2}{6:23}.\begin{aside}[The shelf in the library]
Monetary economics is a field of big books: to get anywhere one has to write a treatise. Mehrling chose it after
seeing the shelf of those books in the library and thinking that a life spent reading them from left to right
would be the best life ever \ts{1}{2}{7:04}.
\end{aside}

% ------------------------------------------------------------------
\section{What is a bank, a shadow bank, a central bank}\label{sec:fp-banks}

The analytical apparatus of the course is, largely, the balance sheet \ts{1}{3}{0:07}. Not balance sheets as
accounting law (what GAAP prescribes), but as a structure that forces rigorous thinking about the economy
\ts{1}{3}{0:27}.\begin{aside}[GAAP]
GAAP, \emph{Generally Accepted Accounting Principles}, are the accounting standards set in the United
States by the Financial Accounting Standards Board. The course ignores them: its balance sheets are a tool of
economic analysis, not of legal reporting.
\end{aside} We introduce three new concepts.

\begin{definition}[New concepts: balance sheet, solvency, liquidity]\label{def:fp-bs}
\leavevmode
\begin{enumerate}[(i)]
  \item A \emph{balance sheet} lists what an agent owns, its \emph{assets} (left side), and what it owes, its
        \emph{liabilities} (right side). \emph{Net worth} (equity) is the balancing item,
        \[ \text{net worth} = \text{assets} - \text{liabilities}, \]
        so that the two sides always have the same total. Drawn as a ``T-account'', the balance sheet has assets
        on the left and liabilities on the right \ts{1}{3}{0:49}.
  \item An agent is \emph{solvent} if the market value of its assets exceeds the market value of its
        liabilities, i.e.\ net worth is positive \ts{1}{3}{2:37}. (From Latin \emph{solvere}, to loosen,
        hence to pay off a debt.)
  \item An agent is \emph{liquid} if it can make the payments it has to make when they come due: it has cash, or
        can obtain it \ts{1}{3}{3:23}. (From Latin \emph{liquidus}, flowing: liquid assets flow easily into
        means of payment.)
\end{enumerate}
\end{definition}

\paragraph{A bank.} A generic bank holds loans, securities (for example Treasury bills) and cash reserves, and
owes deposit accounts and other borrowing (bonds, short-term money market borrowing); net worth balances the two
sides \ts{1}{3}{1:09}. ``Other borrowing'' will be a very big part of the course \ts{1}{3}{1:56}.
\begin{taccts}
\tacct[3cm]{Bank}{Loans\\Securities\\Cash reserves}{Deposits\\Other borrowing\\Net worth}
\end{taccts}
Two issues arise \ts{1}{3}{2:17}. The first is \emph{solvency}: are the assets worth more than the liabilities?
Much of the discussion after 2008, and in the euro crisis, is about this; perhaps entire national banking systems
are insolvent \ts{1}{3}{2:37}. The second, which is what makes banks different from other corporations and which
the course puts first, is \emph{liquidity} \ts{1}{3}{2:58}. Deposits are payable on demand: when a depositor pays
someone at another bank, the bank must deliver cash reserves. Does it have them? If not, can it acquire them, by
borrowing in the money market, by borrowing from the central bank (the Fed's ``discount window''), by selling or
\emph{repo}-ing a security \ts{1}{3}{3:44}? Or must it sell its loans at a \emph{fire-sale} price, below their
value?

\begin{intuition}[Liquidity can become solvency]
A fire sale turns a liquidity problem into a solvency problem: assets sold below their value reduce net worth,
and enough of them can make it negative \ts{1}{3}{4:04}. Most economists approach banking from solvency; the
course approaches it from liquidity, because that is how bankers think: they face the liquidity constraint every
day, when they have to clear their accounts \ts{1}{3}{4:25}. This is why the course starts with the payment system.
A rough analogy: solvency is a statement about stocks (values at one instant), liquidity a statement about flows
(cash in versus cash out, every day). A system can have enough ``energy'' in total and still fail because it
cannot deliver it at the rate required.
\end{intuition}\begin{example}[A fire sale]
A bank has loans 100, reserves 5; deposits 95, net worth 10. Depositors withdraw 20. Reserves cover 5; the bank
must raise 15 by selling loans, and in a panic buyers pay only 60\% of face value: it sells loans of face value
25. Now loans 75, reserves 0, deposits 75, net worth 0. A run of 20 out of 95 wiped out all the capital, although
before it the bank was solvent by 10.
\end{example}

\paragraph{A shadow bank.} Mehrling's research group had just finished a paper on shadow banking, discussed at the
end of the course (lecture 21) \ts{1}{3}{4:45}. A shadow bank, in the sense of the course, is an entity that holds,
for example, residential mortgage-backed securities (loans turned into a kind of bond), strips most of the risk
out of them with an interest rate swap and a credit default swap, and uses what is left, a short-term and
relatively risk-free asset, as collateral to borrow in the \emph{wholesale} money market \ts{1}{3}{5:06}. It
borrows not from households but from money market mutual funds and large corporate investors, through repurchase
agreements (repo), Eurodollar borrowing, asset-backed commercial paper \ts{1}{3}{6:15}.

\begin{definition}[New concepts: shadow bank, wholesale money market]\label{def:fp-shadow}
A \emph{shadow bank} is an entity that does what a bank does (holds long-term credit and finances it with
short-term liabilities that its creditors regard as money-like) outside the regulated deposit-taking banking
system. Its funding comes from the \emph{wholesale money market}, the market where institutions, not households,
lend to each other for short terms, often against collateral. The term ``shadow banking'' was coined by the
economist Paul McCulley (PIMCO) in 2007.
\end{definition}

\begin{taccts}
\tacct[4.2cm]{``Shadow bank''}{Mortgage-backed security\\\quad tranche\\Interest rate swap\\Credit default swap}%
{Money market borrowing:\\\quad repurchase agreements\\\quad asset-backed commercial paper\\\quad auction rate securities\\\quad Eurodollar borrowing}
\end{taccts}
The instruments on this balance sheet are only named here; each gets its own lecture (repo in lecture~7,
Eurodollars in lecture~8, swaps in lectures 19--20). ``It's not rocket science'' \ts{1}{3}{6:36}. The point now is
that this model of credit became dominant: the amount of credit in the US that ran through it became larger than
the amount that ran through traditional banks; the textbooks are all about banks, reality increasingly about
shadow banks \ts{1}{3}{6:56}. Again there is an issue of solvency and an issue of liquidity, but liquidity is
harder: it means being able to \emph{roll over} the short-term funding. If the money market freezes, the shadow
bank cannot roll over and has to liquidate its assets; much of the financial crisis of 2007--2009 was exactly
this \ts{1}{3}{7:38}. The regulatory system was built for traditional banks; the Dodd--Frank Act of 2010 ``doesn't
touch'' shadow banking \ts{1}{3}{8:20}.

\paragraph{The central bank is a bank.} The central bank is also a bank, at a higher level \ts{1}{3}{8:40}. Its
balance sheet is published every week (the Fed's release H.4.1, ``Factors Affecting Reserve Balances of
Depository Institutions''). The release shown in class, of 30 August 2012, gives, in round numbers
\ts{1}{3}{9:00}\ts{1}{3}{11:04}:
\begin{taccts}
\tacct[5.6cm]{Federal Reserve, 30 Aug 2012 (\$ billion)}{Treasuries (no bills)\amt{$\approx$1{,}600}\\Agency debt (Fannie Mae, Freddie Mac)\\Mortgage-backed securities\amt{852}\\Maiden Lane (AIG, Bear Stearns)\\Central bank liquidity swaps\\\textbf{Total}\amt{$\approx$2{,}800}}%
{Currency in circulation\amt{$\approx$1{,}000}\\Reserve balances of banks\amt{$\approx$1{,}500}\\Other\\\\\\\textbf{Total}\amt{$\approx$2{,}800}}
\end{taccts}
\begin{itemize}
  \item The Fed had sold all its short-term Treasury bills and held only notes and bonds: it was moving into
        long-term debt \ts{1}{3}{9:00}.
  \item It held mortgage-backed securities, \$852 billion of them, for the first time in its history: a legacy of
        the crisis \ts{1}{3}{9:21}\ts{1}{3}{9:42}. The ``Maiden Lane'' entries are what was left of the rescues of
        Bear Stearns and AIG.
  \item \emph{Central bank liquidity swaps} are loans of dollars from the Fed to other central banks: the ECB, the
        Swiss National Bank, the Bank of Japan \ts{1}{3}{10:03}.
  \item The total was about \$2.8 trillion, against less than \$1 trillion before the crisis: the balance sheet
        roughly tripled \ts{1}{3}{11:04}.
  \item On the liability side, currency (``green pieces of paper'') was about \$1 trillion, and the reserve
        balances that banks hold as deposits at the Fed were \$1.5 trillion, against about \$50 billion before the
        crisis (as said in class), an order of magnitude change \ts{1}{3}{11:24}\ts{1}{3}{11:45}.
\end{itemize}
This transformation of the Fed's balance sheet is what ``QE1, QE2, QE3'' (quantitative easing) refer to.

\begin{intuition}[The central bank does not ``print money'']
Textbooks write money demand $=$ money supply, $M^d=M^s$, as if the central bank decided a quantity. Behind the
symbols are balance sheets: the central bank is a bank with assets and liabilities, and when it ``expands the money
supply'' it expands \emph{both sides} of its balance sheet at once, for example buying a security (asset) and paying
for it by crediting a bank's reserve account (liability) \ts{1}{3}{10:23}\ts{1}{3}{10:44}. This is a classic banking
operation, and it happens at every level of the system, not only at the top.
\end{intuition}

\begin{taccts}
\tacct[3.2cm]{Fed}{$+$Security}{$+$Reserves of Bank A}\qquad
\tacct{Bank A}{$-$Security\\$+$Reserves}{}
\end{taccts}

% ------------------------------------------------------------------
\section{Central themes}\label{sec:fp-themes}

\subsection{The plan of the course}
The course is about banking \emph{phenomena}, not about banks as legally defined institutions (Mehrling's notes).
It is organised around four functions of banking \ts{1}{4}{0:08}\ts{1}{4}{1:29}, summarised in
\cref{fig:fp-map}:
\begin{enumerate}
  \item \textbf{Banking as a clearing system}: the payment system. Understanding payments is the best starting point
        to understand \emph{liquid assets}.
  \item \textbf{Banking as market making}: once liquid assets are understood, \emph{liquid markets}. These first two
        functions are the subject of the first half of the course, up to the midterm.
  \item \textbf{International money and banking}: the same two ideas extended to the global system \ts{1}{4}{0:28}.
  \item \textbf{Banking as advance clearing}: how finance (forwards, futures, swaps and other derivatives) enters the
        picture \ts{1}{4}{0:48}.
\end{enumerate}
At the end come regulation, what to do about shadow banking, and ``touching the elephant'', the connection of the
money view with what one learns in economics and in finance \ts{1}{4}{1:09}. The fourth classical function,
\emph{intermediation} (channelling savings to investment), is the one textbooks start from; here it comes last.

\subsection{The four prices of money}
In a sense the whole course is about understanding where four prices come from \ts{1}{4}{1:49}.

\begin{definition}[New concepts: the four prices of money]\label{def:fp-four}
\leavevmode
\begin{enumerate}[(i)]
  \item \emph{Par}: the price of one kind of money in terms of another kind of money, today, normally fixed at one
        for one (a deposit of \$107 at a bank exchanges for \$107 of currency) \ts{1}{4}{2:34}\ts{1}{4}{3:16}.
        (Latin \emph{par}, equal.)
  \item \emph{Interest}: the price of money today in terms of money tomorrow; an interest rate $i$ means that $1$
        today exchanges for $1+i$ in one period \ts{1}{4}{2:14}. Its markets: Fed funds, Eurodollars, repo.
        (Medieval Latin \emph{interesse}, ``to be between, to matter'': the compensation due to a lender for the
        loss of the use of the money.)
  \item \emph{Exchange rate}: the price of domestic money in terms of foreign money, or in terms of the world reserve
        money \ts{1}{4}{4:58}.
  \item \emph{Price level}: the price of money in terms of commodities, i.e.\ the inverse of the purchasing power of
        money \ts{1}{4}{5:59}.
\end{enumerate}
\end{definition}

Economics focuses on one of them, the interest rate, the price of money in the sense of a deposit transforming money
today into money tomorrow \ts{1}{4}{2:14}. The other three deserve more attention than they usually get.

\paragraph{Par, the forgotten price.} Par is a more primitive price than interest, and it is neglected
\ts{1}{4}{2:34}. In the bank balance sheet of \cref{sec:fp-banks} deposits (private money, a promise of the bank) sit
on one side and cash reserves (central bank money) on the other, and we take for granted that they trade at par
\ts{1}{4}{2:55}. In financial crises par is sometimes broken: bank money trades at a different rate from central bank
money \ts{1}{4}{3:16}. The system is \emph{hybrid}, part private money and part state money, and par masks the
hybridity: it makes the two look like the same thing \ts{1}{4}{3:36}. When the system works well they \emph{are} the
same; but why does it work well, how do we make it work well, and what do we do in a crisis \ts{1}{4}{3:56}? Par is
key to the payment system. If a Californian seller charged a buyer in New York an extra 5\% because there is no par
clearing between the two states, we would be outraged; yet this is how it was in the United States before the Fed
\ts{1}{4}{4:17}.

\begin{intuition}[Par is a fixed price, and fixed prices are hard]
Par clearing had to be created by institutions, and it is not easy: a price fixed at one for one, which must not
budge whatever the pressure, is like rent control, and economists know how hard such prices are to hold
\ts{1}{4}{4:37}. The rest of the course explains what holds it: the central bank, the clearing system, the
dealers who stand ready to trade at that price.
\end{intuition}\begin{history}[Exchange charges before the Fed]
Before the Federal Reserve, checks drawn on distant banks were often paid at less than face value: the paying bank
deducted an ``exchange charge'' to cover the cost of shipping currency. Universal par clearing was one of the aims of
the Federal Reserve Act (1913); the Fed's campaign for it lasted into the 1920s, and some small banks kept charging
non-par fees for decades.\par
\emph{References:} J.~Jessup, ``The Federal Reserve and Par Clearing'', Federal Reserve Bank of St.~Louis; Young,
``Mobilizing Banking Credits'' (the reading of this week).
\end{history}

\paragraph{The exchange rate.} It is particularly obscured for people living in the United States, because the dollar
is the world's best money \ts{1}{4}{5:18}. Americans never worry that ``the dollar'' could fall below par with the
\emph{international} dollar, the dollar actually used as world reserve \ts{1}{4}{5:39}. Yet there are two dollars, again
a private one and a state one: the international dollar consists largely of dollar deposits at banks outside the US,
mostly European, that have no account at the Fed. These are the \emph{Eurodollars} (lecture~8).

\paragraph{The price level.} Economists usually treat it with the quantity theory of money, $MV=PY$ (or $MV=PT$),
imagining that the price level moves one for one with the quantity of money \ts{1}{4}{5:59}. The course comes to this
only at the end \ts{1}{4}{6:20}; most of the attention goes to par, interest and the exchange rate.

\begin{finance}[The FT market data page]
Two or three pages from the back of the FT there is a ``horrific'' page of tiny numbers \ts{1}{4}{6:40}: market
interest rates (US dollar LIBOR, US dollar certificates of deposit, US overnight repo, Federal funds effective rate),
LIBOR in pounds and Swiss francs, exchange rates, interest rate swap rates \ts{1}{4}{7:00}. By the end of the course
every number on that page should be understandable, and explainable to a roommate: it is not an arcane world reserved
to partners of investment banks, but something the average citizen can and should understand \ts{1}{4}{7:20}.
\end{finance}

\subsection{Lombard Street, then and now}
The inspiration of the course is Bagehot's \emph{Lombard Street} (1873), and Mehrling's \emph{New Lombard Street} is an
homage to it \ts{1}{4}{7:41}. Bagehot wrote a bestseller to explain the money market to his fellow citizens after a
series of financial crises in which the Bank of England had done what it had promised not to do; he argued that it
had done the right thing: lend freely, at a high rate, against good collateral \ts{1}{4}{8:01}. The 2007--2009 crisis was
a crisis of the shadow banking system, and the Fed too did things it had sworn it would never do \ts{1}{4}{8:21}. To get
out of a global crisis, citizens need to understand how the system works, so that every proposal does not get stuck in
``pro-banker'' versus ``anti-banker'' politics; the course takes no side, and approaches the system as scientists and as
citizens \ts{1}{4}{8:42}\ts{1}{4}{9:02}.

\paragraph{Globalization: back to the 19th century.} The world has changed: it is a global financial world
\ts{1}{4}{9:23}. The world of the 1950s textbooks had hardly any private capital markets, and was not global, because
the global financial markets of the 19th century had broken down \ts{1}{4}{9:43}. In terms of globalization the
present resembles the world before the First World War, when the pound sterling was king and the Bank of England
acted as central bank of the whole world, more than the world after the Second (Mehrling's notes). Hence
\begin{center}
\emph{Bagehot is a better guide for today than Irving Fisher} \ts{1}{4}{10:03}:
\end{center}
Fisher lived in a simpler, closed world; Bagehot lived in a global financial system organised around the pound. He did
not know derivatives or high-frequency trading, but he would know how to start thinking about them, because he
understood how money markets work \ts{1}{4}{10:24}.

\paragraph{Finance: all banking is a swap of IOUs.} The second novelty is the integration of finance. Connecting money
and finance, the subjects of Mehrling's two books, which people regard as separate fields, occupied him for ten years
\ts{1}{4}{10:44}. The motto of the course is the connection \ts{1}{4}{11:04}:

\begin{keyformulas}
\textbf{All banking is a swap of IOUs.} \ts{1}{4}{11:25}

\smallskip
\begin{taccts}
\tacct[2.8cm]{Bank}{$+$Loan\\\quad(borrower's IOU)}{$+$Deposit\\\quad(bank's IOU)}\qquad
\tacct[2.8cm]{Borrower}{$+$Deposit\\\quad(bank's IOU)}{$+$Loan\\\quad(borrower's IOU)}
\end{taccts}
A loan is an exchange of promises: the borrower's promise to pay later against the bank's promise to pay on demand.
\end{keyformulas}

By ``swap'' Mehrling means interest rate swaps and credit default swaps, but also ordinary loans \ts{1}{4}{11:25}.
Finance believes it knows everything about swaps; the money view looks at them differently: every swap has a
\emph{funding} piece and a \emph{liquidity} piece, and the liquidity piece drives the action \ts{1}{4}{11:45}. The
course is about money markets and liquidity, connected to solvency, but starting from where Bagehot would start:
the money market as a global system of clearing \ts{1}{4}{12:06}. The aim is monetary theory for the real world of
financial globalization, using as guideposts the great texts of the past and the ``great texts of the present'', Stigum
and the FT \ts{1}{4}{12:26}\ts{1}{4}{12:47}.\begin{aside}[IOU]
\emph{IOU} is a phonetic abbreviation of ``I owe you'', used since the 17th century for an informal written
acknowledgement of a debt. In the course it means any promise to pay, from a handwritten note to a Treasury bond or a
deposit.
\end{aside}

\begin{nfigure}
\begin{tikzpicture}[font=\small,
  box/.style={draw,rounded corners=2pt,align=center,inner sep=3pt,text width=2.6cm,minimum height=1.2cm},
  arr/.style={-{Stealth[length=4pt]},semithick}]
  \node[box,draw=defcol,fill=defcol!8] (c) at (0,0) {\textbf{1. Clearing}\\\scriptsize payment system,\\liquid assets};
  \node[box,draw=defcol,fill=defcol!8] (m) at (3.5,0) {\textbf{2. Market making}\\\scriptsize dealers,\\liquid markets};
  \node[box,draw=intcol,fill=intcol!8] (i) at (7,0) {\textbf{3. International}\\\scriptsize the same, global:\\FX, key currencies};
  \node[box,draw=fincol,fill=fincol!8] (a) at (10.5,0) {\textbf{4. Advance clearing}\\\scriptsize forwards, futures,\\swaps};
  \draw[arr] (c)--(m); \draw[arr] (m)--(i); \draw[arr] (i)--(a);
  \node[draw=srccol,fill=srccol!8,rounded corners=2pt,align=center,inner sep=3pt] (e) at (8.7,-2.1)
    {\textbf{Shadow banking, regulation}\\\scriptsize ``touching the elephant'':\\\scriptsize money view vs economics vs finance};
  \draw[arr] (a.south) -- (a.south |- e.north);
  \draw[decorate,decoration={brace,amplitude=5pt,mirror},thick] (-1.4,-0.8) -- (4.9,-0.8)
    node[midway,below=6pt,font=\scriptsize]{before the midterm (lectures 1--12)};
  \node[draw=keycol,fill=keycol!8,rounded corners=2pt,align=center,inner sep=3pt] (p) at (3.5,1.9)
    {\textbf{four prices of money}\\\scriptsize par $\cdot$ interest $\cdot$ exchange rate $\cdot$ price level};
  \draw[arr,keycol,dashed] (p) -- (c.north); \draw[arr,keycol,dashed] (p) -- (m.north); \draw[arr,keycol,dashed] (p) -- (i.north);
  \node[draw=intucol,fill=intucol!8,rounded corners=2pt,align=center,inner sep=3pt] (s) at (10.5,1.9)
    {\textbf{all banking is}\\\textbf{a swap of IOUs}};
  \draw[arr,intucol,dashed] (s) -- (a);
\end{tikzpicture}
\caption{Map of the course \ts{1}{4}{0:08}: the four functions of banking in the order in which they are treated,
the four prices that the course explains, and the motto that connects money and finance.}\label{fig:fp-map}
\end{nfigure}

% ------------------------------------------------------------------
\section{Reading: Allyn Young}\label{sec:fp-young}

The first reading is by Allyn Young, one of the authors who started Mehrling on his path: the first three chapters of
\emph{The Money Interest and the Public Interest} are a biography of Young \ts{1}{5}{0:00}.

\begin{reading}[Allyn Young, four encyclopedia chapters (1924)]
Young wrote the chapters, unsigned, for an encyclopedia, \emph{The Book of Popular Science}, in 1923--1924, just after
the First World War \ts{1}{5}{0:00}: ``The Mystery of Money'', ``The Monetary System of the United States'',
``Mobilizing Banking Credits'', ``Dear and Cheap Money''. The Fed had been founded in 1913, and the United States had
no tradition of thinking about central banking; Americans were traditionally suspicious of it \ts{1}{5}{0:40}. Young
writes about the Fed when it is a new institution whose role is unclear: he is a friend of the president of the
Federal Reserve Bank of New York and advises him, while open market operations and stabilisation policy are being
invented \ts{1}{5}{1:15}. He compares the Fed with the Bank of England and predicts that London will remain the centre
of world finance; he was wrong: it was the end of sterling's hegemony, although the dollar became the world reserve
currency only after the Second World War, after the Great Depression \ts{1}{5}{1:15}\ts{1}{5}{1:48}. The chapters are
easy to read but deep: on a second reading one notices how carefully Young steers away from minefields
\ts{1}{5}{2:19}. Their spirit is the spirit of the course: having lived through a crisis we do not fully understand,
in a system that is changing, we have to figure it out, as he did in 1924 \ts{1}{5}{1:48}.

\smallskip
\emph{Study questions} (\file{Allyn Young(study).pdf}): where does Young draw the line between money and credit?
Represent with T-accounts the three ways to obtain a bank deposit (\cref{ex:fp-deposits}). How is Young's hostility
to government ``fiat'' money and his sympathy for bank ``credit money'' consistent with the hierarchy of money
(lecture~2)? The reading is analysed in detail in lecture~3.
\end{reading}\begin{history}[Allyn Young (1876--1929)]
Allyn Abbott Young taught at Cornell, Stanford and Harvard; in 1927 he moved to the London School of Economics, where
he died of pneumonia in March 1929, during an influenza epidemic, at 52. He is best known today for ``Increasing
Returns and Economic Progress'' (\emph{Economic Journal}, 1928). The friend at the New York Fed was Benjamin Strong,
its first governor (1914--1928), the dominant figure of the early Fed. Young's encyclopedia chapters are reprinted in
P.~Mehrling and R.~Sandilands (eds.), \emph{Money and Growth: Selected Papers of Allyn Abbott Young}, Routledge, 1999,
from which the course reading is taken.\par
\emph{References:} Mehrling, \emph{The Money Interest and the Public Interest}, 1997, ch.~1--3; Mehrling and
Sandilands, 1999.
\end{history}

% ------------------------------------------------------------------
\flushside
\section*{Key formulas and ideas}
\begin{keyformulas}
\begin{tabularx}{\linewidth}{@{}l L@{}}
Balance sheet & assets $=$ liabilities $+$ net worth (always balances)\\
Solvency & net worth $=$ market value of assets $-$ market value of liabilities $>0$ (a stock)\\
Liquidity & cash available or obtainable $\ge$ payments due, every day (a flow)\\
Four prices of money & par (money vs money, today); interest (money today vs money tomorrow, $1\mapsto 1+i$);
exchange rate (domestic vs foreign money); price level (money vs goods)\\
Fisher diagram & $c_1 + c_2/(1+r) = W$, slope $-(1+r)$; optimum: $u_1/u_2 = 1+r$ \eqref{eq:fp-budget}\\
Central bank & expanding the money supply $=$ expanding both sides of its balance sheet\\
Motto & all banking is a swap of IOUs\\
\end{tabularx}
\end{keyformulas}

\section*{Exercises}

\begin{exercise}[Three ways to get a deposit]\label{ex:fp-deposits}
Using T-accounts for you, your bank and (where needed) a second person and a second bank, record:
\begin{enumerate}[(a)]
  \item you deposit \$100 of currency;
  \item you deposit a check for \$100 written by someone with an account at another bank;
  \item the bank grants you a loan of \$100 by crediting your deposit account.
\end{enumerate}
In which case does the total quantity of deposits in the economy increase? (From the study questions on Young.)
\end{exercise}

\begin{exercise}[Liquidity and solvency]
A bank has loans 200, securities 50, reserves 10; deposits 230, net worth 30. Depositors withdraw 40.
\begin{enumerate}[(a)]
  \item The bank sells securities at their market value to pay. Write the new balance sheet. Is the bank solvent?
        Liquid?
  \item The securities can only be sold at 80\% of their value and the loans at 50\%. Which assets does the bank sell
        first, and what is its net worth afterwards?
  \item Name two ways the bank could meet the withdrawal without selling any asset.
\end{enumerate}
\end{exercise}

\begin{exercise}[Which price of money?]
For each of the following, say which of the four prices of money is involved:
\begin{enumerate}[(a)]
  \item a money market fund announces that each share is now worth \$0.97 instead of \$1;
  \item the Fed funds effective rate is 0.15\%;
  \item a tourist pays 1.30 dollars for one euro;
  \item the consumer price index rises by 2\% in a year;
  \item before 1914 a New York bank paid a check drawn on a Chicago bank at \$99.85 per \$100.
\end{enumerate}
\end{exercise}

\begin{exercise}[The Fisher diagram]
An agent has endowment $(y_1,y_2)=(100,0)$ and utility $u=\ln c_1+\beta\ln c_2$.
\begin{enumerate}[(a)]
  \item Find the optimal $(c_1,c_2)$ as a function of $r$ and $\beta$, and the amount lent.
  \item A second agent has endowment $(0,110)$ and the same utility. For $\beta=1$, find the real rate $r$ that clears
        the market for loans between the two.
\end{enumerate}
\end{exercise}

\begin{exercise}[Quantitative easing in T-accounts]
The Fed buys \$100 billion of mortgage-backed securities from a primary dealer, which holds its account at Bank B.
Write the changes on the balance sheets of the Fed, Bank B and the dealer. Which item of the Fed's liabilities
grows? Relate your answer to the reserve balances of \cref{sec:fp-banks}.
\end{exercise}

\begin{exercise}[A shadow bank in a freeze]
A shadow bank holds mortgage-backed securities worth 100, financed by overnight repo of 95 and capital of 5. One
morning lenders agree to roll over only 80.
\begin{enumerate}[(a)]
  \item How much must the shadow bank raise, and from what?
  \item If buyers pay only 90\% of value for the securities, is the shadow bank still solvent after the sale?
  \item Compare with the bank of \cref{def:fp-bs}: which liability of a traditional bank plays the role of the repo?
\end{enumerate}
\end{exercise}
